\documentclass{testudomanual}

% Compile-time fallback for non-macOS build environments only.
% On macOS the class keeps using Avenir Next / Helvetica Neue / Menlo.
\IfFontExistsTF{Helvetica Neue}{}{%
    \setmainfont{DejaVu Sans}%
    \setsansfont{DejaVu Sans}%
    \setmonofont{DejaVu Sans Mono}[Scale=0.82]%
}

\newcommand{\ManualVersion}{1.7}
\newcommand{\DocumentedTestudoVersion}{0.1.11}
\newcommand{\DocumentedTestudoBuild}{12}
\newcommand{\ManualDate}{October 2026}
\newcommand{\TestudoIconPath}{../../Sources/Testudo/Resources/testudo_icon_blue.png}

\title{Testudo User Manual}
\author{Foivos Karakostas}
\date{\ManualDate}

\begin{document}

% ============================================================
% TITLE PAGE
% ============================================================
\begin{titlepage}
\centering
\vspace*{16mm}
\includegraphics[width=0.22\textwidth]{\TestudoIconPath}
\vspace{13mm}

{\Huge\bfseries\color{TestudoDarkBlue}Testudo\par}
\vspace{4mm}
{\LARGE User Manual\par}
\vspace{7mm}
{\large A native, local-first macOS environment for structured professional and research work\par}

\vfill
\begin{tabular}{rl}
Manual version: & \ManualVersion \\
Application version: & \DocumentedTestudoVersion\ (build \DocumentedTestudoBuild) \\
Platform: & macOS 26+, Apple Silicon \\
Date: & \ManualDate
\end{tabular}

\vfill
{\small Foivos Karakostas\par}
\vspace{8mm}
\end{titlepage}

\pagenumbering{roman}

\chapter*{About this manual}
\addcontentsline{toc}{chapter}{About this manual}

This manual documents the implemented behaviour of \Testudo\ \DocumentedTestudoVersion\ (build \DocumentedTestudoBuild). It is intentionally more than a feature list. It is written so that a new user can understand the data model, perform ordinary work, and actively verify that the application behaves as expected.

Each major chapter therefore contains practical procedures and \term{Try it} exercises. These exercises are safe to perform in the Demo Environment or in a disposable Work Environment. They are especially useful when learning the application after an update.

\manualnote{Version covered}{
This edition is aligned with Testudo 0.1.11. The release introduces deadline-driven Task ordering in the middle column and recursive Sub-task hierarchy, changes Task \menu{Contains} to a newest-first chronological sequence, makes Organization and Group Related Work explicit to the selected entity, improves structural Notes editing, and replaces the generic Work relationship Entity dropdown with Testudo's native People, Groups and Organizations selector. Existing \fileext{.testudoenv} data remains compatible because these changes do not require a new storage representation.
}

\manualwarning{Early-release software}{
Testudo is an early public release. Keep independent backups of important \fileext{.testudoenv} packages. When experimenting with unfamiliar operations, prefer the Demo Environment first.
}

\manualnote{Terminology}{
An \term{Activity} is a chronological Work item that records something that happened. A \term{Calendar Event} is a separate scheduled object. The two concepts are deliberately different even when they refer to the same real-world work.
}

\clearpage
\tableofcontents
\clearpage
\pagenumbering{arabic}

% ============================================================
% CHAPTER 1
% ============================================================
\chapter{What Testudo Is}

\section{Purpose}

Testudo is a native macOS application for work that is difficult to represent with a flat task list. It combines actionable work, durable context, a chronological work log, professional entities, semantic classification, calendars, and local portable storage.

Typical use cases include:
\begin{itemize}
    \item scientific research and laboratory work;
    \item research software engineering;
    \item technical and infrastructure operations;
    \item academic collaborations and committees;
    \item long-running projects with many sub-problems;
    \item work involving multiple institutions, teams or people;
    \item activity where preserving the history of decisions and outcomes is important.
\end{itemize}

The core design principle is that \term{structure} and \term{chronology} are both first-class. A Task can have children and Themes, but Testudo also preserves when the work started, completed, was discontinued, was closed, or was reopened.

\section{The main domains}

Testudo connects several domains that are often kept in separate applications:

\begin{longtable}{L{0.25\textwidth} L{0.65\textwidth}}
\toprule
\textbf{Domain} & \textbf{Purpose} \\
\midrule
\endhead
Work & Tasks, Notes and Activities, including unlimited Task nesting. \\
Themes & Hierarchical semantic classification, independent from Task hierarchy. \\
Entities & People, Groups and Organizations describing professional context. \\
Relationships & Semantic links such as For, Requested by, With, Assigned to and Related to. \\
Calendar & Testudo events plus read-only external calendar sources. \\
History & Significant lifecycle and relationship events. \\
Work Environments & Portable, independent \fileext{.testudoenv} packages. \\
Identity & A local Testudo profile plus separate Environment memberships. \\
\bottomrule
\end{longtable}

\section{Local-first architecture}

Testudo does not require a Testudo cloud service. Work data lives in the selected \fileext{.testudoenv} package. The application also stores installation-level information such as the local user profile, registered Environment locations, filesystem bookmarks and remembered sign-ins.

This separation matters:
\begin{itemize}
    \item copying a \fileext{.testudoenv} package copies the Work Environment;
    \item signing out from the local Testudo installation does not delete the Environment package;
    \item forgetting an Environment from one Mac does not delete its package;
    \item a cloud-synchronised folder can transport the package, but Testudo is still operating on files rather than a central multi-user server.
\end{itemize}

\manualtip{Try it: identify the three layers}{
After opening Testudo, identify: (1) your local Testudo user, (2) the active Work Environment, and (3) your Environment membership. These are related but separate identities.
}

% ============================================================
% CHAPTER 2
% ============================================================
\chapter{Requirements, Installation and Updating}

\section{Supported baseline}

The current source targets:
\begin{itemize}
    \item macOS 26 or later;
    \item Apple Silicon;
    \item an Apple developer toolchain with Swift 6 or later;
    \item Git.
\end{itemize}

The repository uses a source-first installation model. The application is compiled on the Mac where it is used and is ad-hoc signed by the local build script.

\section{Quick installation}

Run:
\begin{lstlisting}
curl -fsSL \
  https://raw.githubusercontent.com/foivoskar/Testudo/main/install.sh \
  | bash
\end{lstlisting}

The installer checks the machine, clones or updates the repository, refuses to overwrite a checkout with local modifications, builds Testudo, creates a DMG, verifies it and opens the finished installer.

When the DMG opens, drag \filepath{Testudo.app} into \filepath{/Applications}.

\section{Compatibility check only}

To test the toolchain without cloning or building:
\begin{lstlisting}
curl -fsSL \
  https://raw.githubusercontent.com/foivoskar/Testudo/main/install.sh \
  | bash -s -- --check
\end{lstlisting}

\section{Manual source build}

\begin{lstlisting}
git clone https://github.com/foivoskar/Testudo.git
cd Testudo
./Scripts/build-local-dmg.sh
\end{lstlisting}

The resulting DMG is written under \filepath{dist/}. For this release the expected name is normally:
\begin{lstlisting}
dist/Testudo-0.1.7-macOS-arm64.dmg
\end{lstlisting}

\section{Updating}

If the source checkout is clean:
\begin{lstlisting}
cd ~/Testudo
git pull --ff-only
./Scripts/build-local-dmg.sh
\end{lstlisting}

Replace the previous copy in \filepath{/Applications} with the newly built application.

\manualtip{Try it: verify the installed version}{
Build and install Testudo, then confirm that the application identifies itself as \DocumentedTestudoVersion\ build \DocumentedTestudoBuild. If you maintain a development checkout, also inspect \filepath{Scripts/version.sh} and confirm that it contains the same version and build number.
}

% ============================================================
% CHAPTER 3
% ============================================================
\chapter{First Launch and the Local Testudo Profile}

\section{Create or restore a profile}

On first use, Testudo presents \menu{Create or Restore Your Profile}. The local profile belongs to this installation, not to a Work Environment.

The setup interface can record professional information including:
\begin{itemize}
    \item first, middle, last and preferred name;
    \item academic title and job title;
    \item professional email, phone and office;
    \item website, ORCID, LinkedIn and GitHub;
    \item professional address, city, postal code and country;
    \item professional fields, responsibilities and notes.
\end{itemize}

Only enter fields that are useful to you. The profile can later be edited from the application settings.

\section{Importing a portable profile}

If you previously exported a \fileext{.testudouser} file, you can import it instead of re-entering profile data.

A \fileext{.testudouser} file restores profile information and avatar data. It does \emph{not} restore:
\begin{itemize}
    \item the application password;
    \item registered Work Environment locations;
    \item Environment passwords;
    \item remembered Environment sign-ins.
\end{itemize}

\section{Application password and recovery}

After creating the profile, Testudo can protect the installation with an application password. The setup explicitly states that this password is separate from every Work Environment password.

The application password must contain at least eight characters. Recovery uses three user-defined security questions.

Security-question answers are stored as salted hashes. Matching normalises capitalization, accents and repeated spaces.

\manualtip{Try it: profile and lock flow}{
Create a disposable local profile, configure the application password and three recovery questions, use \menu{Log Out of Testudo}, and log back in. Confirm that your registered Environment list is still present. Then compare this with \menu{Sign Out from Testudo...}, which is intentionally much more destructive to local installation state.
}

% ============================================================
% CHAPTER 4
% ============================================================
\chapter{Local Log Out, Local Sign Out and Environment Sign Out}

Testudo exposes several actions that sound similar but have different scopes.

\section{Log Out of Testudo}

Logging out locks the application. Your local profile, registered Work Environments and remembered state remain available after the next login.

Use this when you simply want to protect the application while keeping the current installation configured.

\section{Sign Out from Testudo}

This is an installation-level reset for the local Testudo identity. The confirmation screen requires typing \texttt{SIGN OUT}.

It removes local information from this Mac, including the local profile, application password, registered Work Environments, local bookmarks, Environment identity mappings and remembered sign-ins.

It does \textbf{not} delete or modify the actual \fileext{.testudoenv} packages.

Before signing out, Testudo offers \button{Export My Profile Before Signing Out...}.

\section{Sign Out of Environment}

This action affects the active Work Environment session rather than the local application account. It is useful when the same local Testudo installation is used to enter an Environment with different credentials.

\manualwarning{Do not confuse the scopes}{
If you only want to lock the application, use \menu{Log Out of Testudo}. If you want to remove local Testudo-user state from the Mac, use \menu{Sign Out from Testudo...}. If you only want to end the current Environment session, use \menu{Sign Out of Environment}.
}

% ============================================================
% CHAPTER 5
% ============================================================
\chapter{Work Environments}

\section{The Work Environments screen}

After local authentication, Testudo shows \menu{Work Environments} with the message \emph{Choose where you want to work}. Registered Environments appear as cards showing their name and storage path.

Available actions are:
\begin{itemize}
    \item \button{Create New Environment};
    \item \button{Open Existing Environment...};
    \item \button{Load Demo Environment}.
\end{itemize}

\section{Creating an Environment}

The creation sheet asks for:
\begin{itemize}
    \item Environment name;
    \item initial Administrator username;
    \item Administrator password;
    \item password confirmation.
\end{itemize}

The first account belongs to the Work Environment and is separate from the local Testudo profile.

The Environment is saved as a portable \fileext{.testudoenv} package. You can choose a local, cloud-synchronised or external storage location.

Internally, a normal package contains files such as:
\begin{lstlisting}
Example.testudoenv/
|-- EnvironmentManifest.json
|-- EnvironmentData.json
|-- EnvironmentCredentials.json
`-- EnvironmentIcon.png        # optional
\end{lstlisting}

\section{Opening an existing Environment}

Use \button{Open Existing Environment...} and select a \fileext{.testudoenv} package. Testudo validates the package and registers its location with the current installation.

If the Environment requires identity selection, Testudo presents the Environment login flow.

\section{Environment login}

The entry sheet contains a username, password and \menu{Stay signed in} option.

When \menu{Stay signed in} is enabled, this Mac can reopen the Environment without asking for those credentials until the user explicitly signs out of the Environment.

\section{Forgetting an Environment from one Mac}

In Environment settings, removing an Environment from Testudo only forgets its registration on that installation. The \fileext{.testudoenv} package and all Environment data remain at the storage location and can be reopened later.

\manualtip{Try it: create and reopen}{
Create a disposable Environment called \emph{Manual Test}. Close it, return to Work Environments, and reopen it. Then remove its registration from Testudo and use \button{Open Existing Environment...} to register the same package again. Verify that the Work data is still present.
}

% ============================================================
% CHAPTER 6
% ============================================================
\chapter{Environment Membership and Administration}

\section{Environment roles}

Environment memberships currently have two roles:
\begin{itemize}
    \item Administrator;
    \item User.
\end{itemize}

A membership also has an active/inactive state.

\section{Membership identity}

An Environment membership can either be linked to an existing Person entity or maintain its own first and last name. When linked to a Person, the membership name follows the Person.

The management view exposes:
\begin{itemize}
    \item Linked Person;
    \item first and last name when unlinked;
    \item username;
    \item immutable Membership ID;
    \item Environment role;
    \item active membership state.
\end{itemize}

\section{Environment passwords}

Environment passwords are separate from the local application password.

An Administrator can assign a new Environment password without knowing the member's current password. Security-question answers are never revealed to Administrators.

\section{Environment recovery questions}

A member can configure three recovery questions for the current Environment account. These recovery questions belong to that Environment account and travel with the \fileext{.testudoenv} package.

If a member cannot answer a recovery question, an Environment Administrator can still assign a new password.

\manualtip{Try it: verify identity separation}{
In a disposable Environment, create or inspect a second membership. Confirm that changing the Environment password does not change the Testudo application password. If the membership is linked to a Person, edit the Person name and verify that the membership display follows the linked Person.
}

% ============================================================
% CHAPTER 7
% ============================================================
\chapter{The Three-Column Interface}

\section{Sidebar}

The left sidebar is the global navigation layer. The current sidebar contains:
\begin{itemize}
    \item New Entry;
    \item Today;
    \item Calendar;
    \item All Tasks;
    \item To Do;
    \item In Progress;
    \item Completed;
    \item Discontinued;
    \item Archive;
    \item Timeline;
    \item Themes;
    \item Organizations;
    \item Groups;
    \item People.
\end{itemize}

\section{Middle column}

The middle column shows the collection selected in the sidebar: a Task list, Archive tree, calendar, Theme list, entity list or other contextual collection.

Task-oriented middle-column views prioritise Tasks with deadlines. Earlier deadlines appear first, followed by Tasks without deadlines. In \menu{All Tasks}, the same rule applies recursively among Sub-task siblings. Active \menu{Today} work uses the same deadline priority where Task ordering is relevant.

\section{Detail pane}

The right pane shows the selected object's detail view or a creation/editing form. Normal detail views are mostly read-only and expose a single whole-object \button{Edit} action where appropriate.

\section{Sidebar visibility}

Testudo uses the native macOS sidebar behaviour. The system-provided sidebar control hides and restores the first column.

\section{Back and Forward navigation}

The detail pane maintains browser-style navigation history through Back and Forward controls and two-finger trackpad navigation.

The navigation history represents the complete three-column context. When you move backward or forward, Testudo restores:
\begin{itemize}
    \item the sidebar section;
    \item the middle-column selection;
    \item the right-pane detail content.
\end{itemize}

Archived Tasks participate in the same history as active Work.

\manualtip{Try it: verify three-column history}{
Open an active Task, then open \menu{Archive} and select a Closed Task, then open a Person. Use Back and Forward or the trackpad gesture. Verify that each step restores not only the right pane but also the corresponding sidebar section and middle-column selection.
}

% ============================================================
% CHAPTER 8
% ============================================================
\chapter{Work: Tasks, Notes and Activities}

\section{Three Work types}

\begin{description}
    \item[Task] Actionable work with a Task status and optional lifecycle dates.
    \item[Note] Durable context that does not require Task status.
    \item[Activity] A work-log item recording something that happened, optionally at a specific date and time.
\end{description}

Only Tasks can contain child Work. A Task may contain sub-tasks, Notes and Activities, and sub-tasks can contain further Work without a fixed nesting limit.

\section{Creating Work}

Choose \menu{New Entry} or press Command-N. The Work editor includes sections for content, classification, parentage, timing and semantic relationships.

Common fields include:
\begin{itemize}
    \item Type;
    \item one or more Themes;
    \item Parent Task;
    \item optional Title;
    \item Description or Note / What happened? text;
    \item relationships to People, Groups and Organizations.
\end{itemize}

At least a title or description is required.

\section{Editing existing Work}

Open the Work item and choose \button{Edit}. The same complete editor is used for Tasks, Notes and Activities. Existing values are pre-filled.

\section{Creating child Work quickly}

Inside a Task detail page, the \menu{Contains} area can create:
\begin{itemize}
    \item Sub-task;
    \item Note;
    \item Activity.
\end{itemize}

Activities created this way are still Work-log objects, not Calendar Events.

The \menu{Contains} list is shown newest first. Activities are positioned by their \term{Occurred} time when one is available; Tasks and Notes use their Created time. Deadline does not control this sequence because \menu{Contains} represents chronology rather than priority.

\manualtip{Try it: build a small hierarchy}{
Create a top-level Task called \emph{Prepare field campaign}. Add a sub-task called \emph{Check instruments}, a Note with background information and an Activity recording a completed phone call. Then add a deeper sub-task under \emph{Check instruments}. Navigate up and down the hierarchy and confirm that each item retains its own detail page.
}

% ============================================================
% CHAPTER 9
% ============================================================
\chapter{Task Workflow Status}

\section{The four statuses}

In Testudo 0.1.7, the user-facing Task statuses are exactly:
\begin{itemize}
    \item To Do;
    \item In Progress;
    \item Completed;
    \item Discontinued.
\end{itemize}

\term{Closed} is deliberately absent from this list because it is an archive lifecycle action rather than a workflow outcome.

\section{To Do}

Use To Do for work that exists but has not begun. A Task can still have a deadline while remaining To Do.

\section{In Progress}

Use In Progress when work has begun. When a Task changes from To Do to In Progress and does not already have a Started date, Testudo can record the transition time automatically. A manually entered Started time remains independently editable.

\section{Completed}

Use Completed when the work ended successfully. The completion timestamp is preserved independently from Created and Started.

\section{Changing status from the detail view}

The status control allows normal workflow changes. Discontinued requires additional metadata, so it is intentionally handled through the full Work editor rather than a shortcut that could omit the reason.

\manualtip{Try it: normal lifecycle}{
Create a Task as To Do. Change it to In Progress and check its Started value. Then mark it Completed and confirm that a completion time appears. Reopen the editor and verify that Created, Started and Completed represent different concepts and can describe historical work accurately.
}

% ============================================================
% CHAPTER 10
% ============================================================
\chapter{Discontinued Tasks}

\section{Why Discontinued exists}

Not every Task finishes successfully. Work may fail, become unnecessary, be replaced by another approach, or prove infeasible. Deleting such a Task loses useful history; marking it Completed is misleading. Testudo therefore provides a terminal \term{Discontinued} status.

\section{Required metadata}

When the Task status is Discontinued, the editor requires a reason. Available reasons are:
\begin{itemize}
    \item Failed;
    \item Superseded;
    \item No longer needed;
    \item Not feasible;
    \item Other.
\end{itemize}

The editor also includes:
\begin{itemize}
    \item Discontinued date/time with source time zone;
    \item optional Outcome / note.
\end{itemize}

The discontinued time cannot be in the future and cannot be earlier than Started when a Started time exists.

\section{What remains visible}

The normal detail page shows the Discontinued time, reason and optional outcome. The Task also appears in the \menu{Discontinued} sidebar section while it remains active rather than archived.

A discontinuation lifecycle event is retained in History.

\section{Historical cycles}

If a Task later returns to active work, earlier discontinuation history remains preserved. A later discontinuation can therefore represent a new lifecycle cycle without erasing the previous one.

\manualtip{Try it: discontinue for a reason}{
Create a Task called \emph{Use vendor API A}. Start it, edit it to Discontinued, choose \emph{Superseded}, and write \emph{Replaced by API B}. Confirm that the Task appears in \menu{Discontinued}, that its detail page shows the reason and outcome, and that its Log contains a discontinuation event.
}

% ============================================================
% CHAPTER 11
% ============================================================
\chapter{Closing, Archive and Reopening}

\section{Closed is not a status}

Closing is an archival lifecycle operation. Only a \textbf{top-level Task} can be Closed, and it must already be terminal:
\begin{itemize}
    \item Completed; or
    \item Discontinued.
\end{itemize}

A To Do or In Progress Task cannot be closed directly.

\section{What closing does}

Closing a top-level Task:
\begin{itemize}
    \item records a Closed timestamp and source time zone;
    \item retains the real workflow status;
    \item moves the whole Task subtree to Archive;
    \item records a Closed lifecycle History event;
    \item removes the subtree from active Work queries.
\end{itemize}

Children do not independently become Closed. They are archived because their top-level ancestor is closed.

\section{Archive presentation}

The Archive middle column displays Closed top-level Tasks and their retained subtree. The root still indicates whether the actual outcome was Completed or Discontinued, and its context shows the Closed time.

\section{Reopen}

A Closed top-level Task exposes \button{Reopen}. Reopening removes the archive marker, returns the tree to active views, preserves the actual Completed or Discontinued status and records a Reopened History event.

\section{Backward compatibility}

Testudo 0.1.6 stored Closed as a Task status. Testudo 0.1.7 keeps a legacy persistence-only value so older Environment data can be decoded, then migrates the item to the separate archive lifecycle while preserving the most meaningful previous workflow state available.

\manualtip{Try it: complete archive lifecycle}{
Perform the following in a disposable Environment:
\begin{enumerate}
    \item Create a top-level Task and a child Note.
    \item Mark the Task Completed.
    \item Use \button{Close} in the detail header.
    \item Confirm that the Task disappears from active lists and appears in \menu{Archive} with its child Note.
    \item Open the archived Task and verify that its Status is still Completed.
    \item Use Back/Forward navigation and confirm Archive selection is restored correctly.
    \item Choose \button{Reopen}.
    \item Confirm that the Task returns to active views as Completed and that Closed/Reopened events are preserved in its Log.
\end{enumerate}
Repeat the exercise with a Discontinued Task and verify that its discontinuation reason and outcome survive Close and Reopen.
}

% ============================================================
% CHAPTER 12
% ============================================================
\chapter{Notes}

Notes are informational Work items. They do not have Task status and are useful for durable context that belongs near the work.

A Note can have:
\begin{itemize}
    \item title and body;
    \item one or more Themes;
    \item a Parent Task;
    \item a Reminder;
    \item semantic relationships.
\end{itemize}

Use a Note when the item should remain available as context but does not represent work that progresses through To Do, In Progress, Completed or Discontinued.

\manualtip{Try it: reminder-bearing Note}{
Inside a Task, create a Note containing a decision or reference, enable a Reminder for tomorrow, and save it. Confirm that it appears under the parent Task and that the reminder is represented in Calendar and relevant Today views when its date becomes current.
}

% ============================================================
% CHAPTER 13
% ============================================================
\chapter{Activities and the Work Log}

\section{Activity semantics}

An Activity records something that happened while working. Examples include a discussion, measurement, intervention, review, decision or resolved problem.

The Activity editor provides an optional \menu{Occurred} date/time with source time zone.

\manualnote{Activity vs Calendar Event}{
An Activity belongs to the Work history. A Calendar Event belongs to the Calendar domain. A meeting that is planned for tomorrow is a Calendar Event; after it occurs, you may also record an Activity containing the useful outcome.
}

\section{Activity in Timeline and History}

Activities contribute to chronological views and appear in the Work Log. A parent Task can therefore preserve a running record without rewriting its description repeatedly.

\manualtip{Try it: scheduled event and resulting Activity}{
Create a Calendar Event for a meeting and link it to a Task. After the simulated meeting, create an Activity under the Task called \emph{Meeting outcome} and record a historical Occurred time. Compare how the Calendar Event and Activity appear in their respective views.
}

% ============================================================
% CHAPTER 14
% ============================================================
\chapter{Themes and Multi-Theme Classification}

\section{Purpose}

Themes classify Work semantically and independently from the Task hierarchy and professional entity hierarchy.

A Theme can have a parent Theme, allowing arbitrary Theme depth.

\section{Theme detail fields}

The current Theme interface exposes information including:
\begin{itemize}
    \item Name and Parent Theme;
    \item Summary and Notes;
    \item Start and Target dates;
    \item Tags;
    \item URL;
    \item SF Symbol;
    \item related People, Groups and Organizations;
    \item Created and Updated timestamps.
\end{itemize}

The normal detail page is read-only and exposes a whole-object \button{Edit} action.

\section{Multiple Themes on one Work item}

A Work item can belong to several Themes simultaneously. This avoids duplicating the same Task simply because it belongs to multiple conceptual domains.

\section{Descendant aggregation}

Related Work for a parent Theme includes Work assigned to descendant Themes recursively. This aggregation does not rewrite direct membership.

\manualtip{Try it: multi-Theme aggregation}{
Create Theme A, then child Theme B. Create a Task assigned directly to Theme B and also to a second independent Theme. Open Theme B and Theme A and verify that the Task appears in Related Work for both levels of the A/B hierarchy, while the Task editor still shows its actual direct Theme selections.
}

% ============================================================
% CHAPTER 15
% ============================================================
\chapter{People}

People are first-class professional entities with richer profile data than a simple name string.

\section{Person profile fields}

The Person editor can represent:
\begin{itemize}
    \item first, middle, last and preferred name;
    \item academic title and position;
    \item primary and secondary email/phone;
    \item office and assistant/contact;
    \item staff or employee ID;
    \item address and workplace;
    \item website, LinkedIn and GitHub;
    \item ORCID, ResearcherID, Scopus Author ID and Google Scholar;
    \item fields/expertise, responsibilities and tags;
    \item preferred language and time zone;
    \item notes;
    \item Organization and Group affiliations.
\end{itemize}

\section{Affiliations}

A Person can belong to multiple Organizations and Groups. One affiliation can be primary. Some affiliations can be shown as automatic when implied by structure or membership data.

\section{Environment membership link}

A Person may be linked to an Environment membership. The Person detail page then exposes membership identity, role, active/inactive state and whether a local Environment password is configured.

A Person linked to an Environment membership cannot be deleted until the membership is unlinked.

\manualtip{Try it: professional directory}{
Create a Person with a preferred name and ORCID, affiliate the Person with an Organization and a Group, then relate a Task to that Person using \term{With}. Open the Person and confirm that the Task appears in Related Work.
}

% ============================================================
% CHAPTER 16
% ============================================================
\chapter{Organizations and Groups}

\section{Organizations}

Organizations model institutional structure such as universities, institutes, departments, companies or observatories. Organizations can themselves be contained in other Organizations.

The editor can expose fields such as full name, short name, website, email, phone, address, city, postal code, country, tags, symbol, dates and notes. Long-form Notes use a larger, clearly bounded writing surface so that editable text is visually distinct from the surrounding form.

\section{Groups}

Groups model teams, laboratories, committees, collaborations and working groups. A Group may be independent or contained by one or more compatible structural entities according to the current hierarchy rules.

\section{Affiliated People}

Organization and Group editors can manage affiliated People. Their Related Work pane is deliberately explicit: a Work item appears only when it has a relationship to the exact selected Organization or Group. Work is not included merely because it belongs to affiliated People, descendant Groups, descendant Organizations or another structural relation. A relationship copied onto a child Work item is stored on that child and therefore continues to count as a direct relationship.

\manualtip{Try it: build a small structure}{
Create an Organization called \emph{Institute}, a child Organization called \emph{Department}, and a Group called \emph{Lab} inside the Department. Affiliate a Person with the Lab. Relate a Task to the Lab and verify Related Work from the relevant structural views.
}

% ============================================================
% CHAPTER 17
% ============================================================
\chapter{Work Relationships}

Work relationships attach professional context directly to Work items.

Available roles are:
\begin{itemize}
    \item For;
    \item Requested by;
    \item With;
    \item Assigned to;
    \item Related to.
\end{itemize}

Each relationship points to a Person, Group or Organization.

In the Work editor, relationship entities are selected through the same native Testudo selector used elsewhere in the application. The selector separates People, Groups and Organizations and supports the existing searchable entity-selection workflow rather than presenting one generic Entity dropdown.

\section{Inheritance by child items}

A relationship may be marked \menu{Inherited by child items}. Child Work can then inherit relevant context from its parent Task.

When changing the Parent Task in the full editor, Testudo may ask whether to use the new Parent's inherited relationships. Copying them replaces the current related People, Groups and Organizations with the Parent relationships marked for inheritance.

\manualtip{Try it: inherited context}{
Create a parent Task related \term{For} an Organization and enable inheritance. Create a child Task and inspect its inherited context. Then move the child to another Parent Task with a different inherited relationship and test the \menu{Copy Parent Relationships} / \menu{Keep Current Relationships} decision.
}

% ============================================================
% CHAPTER 18
% ============================================================
\chapter{Related Work}

People, Themes, Organizations and Groups share a Related Work presentation.

Related Work can contain Tasks, Notes and Activities. When related Work exists, the detail area can present the Work separately from the object's own inspector information.

Aggregation follows the relevant semantic structure:
\begin{itemize}
    \item parent Themes can include descendant Theme Work;
    \item Organizations and Groups can aggregate from relevant descendants;
    \item duplicate Work is suppressed when several paths resolve to the same item.
\end{itemize}

Rows are navigable, so Related Work is also a bridge between the structural and Work domains.

\manualtip{Try it: navigate through Related Work}{
Open a Person or Theme that has several related Work items. Select a related Task, then use Back. Verify that the original Person/Theme and its corresponding left/middle-column context are restored.
}

% ============================================================
% CHAPTER 19
% ============================================================
\chapter{Today}

\section{Today is a work summary, not only a schedule}

The Today view combines several kinds of current-day relevance. In 0.1.7 the interface includes sections such as:
\begin{itemize}
    \item Needs Attention;
    \item In Progress;
    \item Upcoming;
    \item Done Today;
    \item Discontinued Today;
    \item Today's Log.
\end{itemize}

The underlying Today collection considers active Work whose Created, Logged/Occurred, Started, Completed or Discontinued date is today.

\section{Needs Attention}

Needs Attention highlights urgent active items such as overdue deadlines, work due today and relevant reminders. Closed Task trees are excluded because archived Work is not part of the active workflow.

\section{Upcoming}

Upcoming provides near-future actionable context, such as deadlines and reminders.

\manualtip{Try it: populate Today}{
Create one Task due today, start another Task today, complete a third Task today, discontinue a fourth Task today and create an Activity with today's Occurred time. Open \menu{Today} and identify where each item appears. Then Close one terminal top-level Task and confirm the archived tree no longer contributes to the active Today workflow.
}

% ============================================================
% CHAPTER 20
% ============================================================
\chapter{Task Lists and Timeline}

The sidebar exposes dedicated Task filters:
\begin{itemize}
    \item All Tasks;
    \item To Do;
    \item In Progress;
    \item Completed;
    \item Discontinued.
\end{itemize}

These lists operate on active Work. Archived Task trees are excluded until reopened.

\section{Timeline}

Timeline focuses on chronological Work activity rather than a single workflow state. Activities and lifecycle dates help make long-running work understandable after the fact.

\manualtip{Try it: compare filters}{
Create four Tasks, one in each workflow status. Confirm that each appears in its corresponding sidebar filter. Close the Completed and Discontinued top-level Tasks and verify that they move to Archive and disappear from the active status filters. Reopen them and verify that they return to their retained status filters.
}

% ============================================================
% CHAPTER 21
% ============================================================
\chapter{History and the Work Log}

\section{Audit timestamps}

The History section exposes key timestamps such as Created and Updated, plus type-specific lifecycle values such as Started, Completed, Discontinued or Occurred.

\section{Log events}

The chronological Log can represent:
\begin{itemize}
    \item created;
    \item edited;
    \item status changed;
    \item scheduled;
    \item started;
    \item completed;
    \item discontinued;
    \item closed;
    \item reopened;
    \item activity logged;
    \item relationship added or removed;
    \item moved.
\end{itemize}

The Log is deliberately not only a mirror of the current state. Earlier lifecycle cycles remain meaningful historical records.

\section{Deletion and history}

Deleting a Work item is different from discontinuing or closing it. Deletion permanently removes the item and its descendants, plus associated relationships and history for those deleted Work items.

\manualwarning{Prefer lifecycle actions for historical work}{
If a Task existed and its outcome matters, prefer Completed, Discontinued and possibly Close. Use deletion for data that should genuinely be removed rather than preserved as history.
}

% ============================================================
% CHAPTER 22
% ============================================================
\chapter{Calendar Overview}

\section{What appears in Calendar}

The Calendar combines several scheduling concepts:
\begin{itemize}
    \item Calendar Events;
    \item Task deadlines;
    \item Note reminders.
\end{itemize}

The month view includes navigation for previous/next month and \button{Today}. Day summaries distinguish Events, Deadlines and Reminders.

Toolbar actions include:
\begin{itemize}
    \item Connect iCal Calendar;
    \item Connect Apple / iCloud Calendar;
    \item Sync connected external calendars;
    \item New Calendar Event.
\end{itemize}

\section{Calendar sources}

Current source kinds include:
\begin{itemize}
    \item Testudo local calendar;
    \item iCal Subscription;
    \item Apple Calendar / iCloud.
\end{itemize}

\manualtip{Try it: deadline, reminder and event}{
For the same date, create a Task deadline, a Note reminder and a Calendar Event. Open that date in Calendar and confirm that all three appear but are identified by their different scheduling roles.
}

% ============================================================
% CHAPTER 23
% ============================================================
\chapter{Creating and Editing Calendar Events}

\section{New Calendar Event}

A local Calendar Event can contain:
\begin{itemize}
    \item Title;
    \item Calendar;
    \item All day flag;
    \item Start and End;
    \item source time zones for timed events;
    \item Location;
    \item Notes;
    \item Related Tasks;
    \item Related Themes.
\end{itemize}

The End cannot be earlier than Start.

Related Tasks are selected from the Task hierarchy, and related Themes are selected from the Theme hierarchy.

\section{Calendar Event status}

Calendar event statuses include Confirmed, Tentative and Cancelled where applicable to the event model.

\section{External events}

Events imported from external calendars are read-only for their external details. Testudo relationships can still be edited. The detail view explicitly states when an event is a read-only external calendar event.

\manualtip{Try it: event linking}{
Create a local Event called \emph{Project review}, link it to a Task and a Theme, then open the Event and navigate through those links. Return with Back and verify the Calendar context is restored.
}

% ============================================================
% CHAPTER 24
% ============================================================
\chapter{Apple Calendar Integration}

\section{Connecting Apple Calendar}

Choose \button{Connect Apple / iCloud Calendar}. Testudo requests macOS Calendar access and lists calendars available through Apple Calendar on the Mac.

The connection view states that imported events are read-only in Testudo and that only Testudo relationships can be edited.

The initial synchronisation window covers approximately one year in the past through two years in the future.

\section{Permission and availability}

If macOS permission is not granted, or no compatible calendars are available, the connection view reports the condition. Permission is controlled by macOS privacy settings.

\manualtip{Try it: read-only behaviour}{
Connect a non-critical Apple Calendar, sync it and open one imported event. Confirm that event details are read-only while Related Work / Theme links remain manageable inside Testudo.
}

% ============================================================
% CHAPTER 25
% ============================================================
\chapter{iCal Subscriptions}

Choose \button{Connect iCal Calendar} and provide the secret address in iCal format. An optional Calendar name can be entered.

Events imported from an iCal subscription are read-only.

The secret subscription address is treated as sensitive and is stored in macOS Keychain rather than in ordinary Testudo data.

The Keychain credential is local to the Mac on which it was entered. A shared \fileext{.testudoenv} package stores only the identifier used to locate that local credential; it does not contain the Secret iCal address itself.

When the same Work Environment is opened on another Mac, the Calendar and its imported events may therefore already be present while the local Secret address is absent. Use \button{Restore Secret Address} to provide the credential on that Mac. Restoration keeps the existing Calendar identity and updates the existing imported subscription rather than creating a second Calendar.

A missing local Keychain credential is treated as a machine-local condition and is not written into the shared Work Environment as a synchronization error.

\manualwarning{Treat secret iCal URLs as credentials}{
Anyone with a secret subscription URL may be able to read the calendar feed. Do not place the URL in screenshots, Notes or shared documentation.
}

\manualtip{Try it: subscription lifecycle}{
If you have a disposable iCal feed, connect it, trigger external-calendar synchronisation and inspect an event. Verify that the source behaves as read-only and that the feed URL is not shown as ordinary Work Environment data.
}

% ============================================================
% CHAPTER 26
% ============================================================
\chapter{Time Zones and Historical Dates}

\section{Absolute time and source zone}

Testudo stores absolute timestamps as dates and, where relevant, stores a separate IANA time-zone identifier describing the zone in which the value was entered or recorded.

Time-zone-aware fields include, depending on type:
\begin{itemize}
    \item Created and Updated;
    \item Scheduled;
    \item Deadline;
    \item Reminder;
    \item Started;
    \item Completed;
    \item Discontinued;
    \item Closed;
    \item Activity Occurred;
    \item Calendar Event Start and End.
\end{itemize}

The current device time zone is the normal display context, while the source zone helps preserve meaning across travel and historical entry.

\section{Historical entry}

Started, Discontinued and Activity occurrence times can describe events before the item was entered into Testudo, subject to logical validation rules.

\manualtip{Try it: source-zone preservation}{
Create an Activity with a historical time and explicitly choose another IANA time zone. Save it, then reopen the detail and editor. Verify that the absolute moment and source-zone context remain consistent.
}

% ============================================================
% CHAPTER 27
% ============================================================
\chapter{Demo Environment}

\section{Purpose}

\button{Load Demo Environment} creates a complete fictional data set intended for exploration. It includes People, Organizations, Groups, Themes, Tasks, sub-tasks, Notes, Activities, Calendar Events, affiliations, relationships, chronology, history and archived work.

Dates are generated relative to the date on which the Demo is created, so the environment remains useful over time.

\section{Demo naming}

Multiple Demo Environments can coexist. The generator uses names such as:
\begin{lstlisting}
Demo
Demo-2
Demo-3
\end{lstlisting}

Available numbering can be reused after a Demo is removed.

\section{Suggested tour}

Use the Demo to inspect, in order:
\begin{enumerate}
    \item Today;
    \item one multi-level Task tree;
    \item a Person with affiliations and Related Work;
    \item a Theme with descendant-related Work;
    \item Calendar;
    \item Archive;
    \item browser-style Back/Forward navigation across different object types.
\end{enumerate}

\manualtip{Try it: acceptance-test tour}{
Before upgrading an important real Environment, create a fresh Demo and perform the tour above. It is a low-risk way to detect obvious UI or persistence regressions.
}

% ============================================================
% CHAPTER 28
% ============================================================
\chapter{Storage, Portability and Backups}

\section{Portable Environment package}

The \fileext{.testudoenv} package is the principal portable unit. It contains the Environment's manifest, data, credentials and optional icon.

\section{Cloud-synchronised folders}

An Environment can live in Dropbox, iCloud Drive, OneDrive or another filesystem provider visible to macOS. Testudo still works on the package as a file package.

Avoid simultaneous conflicting edits from several Macs. Wait for file synchronisation to complete before opening the same package elsewhere.

\section{Backups}

File synchronisation is not a substitute for independent backup. A deletion or corruption can itself be synchronised.

Recommended practice:
\begin{enumerate}
    \item keep at least one independent backup of important Environments;
    \item before upgrading between early versions, copy critical \fileext{.testudoenv} packages;
    \item avoid manually editing JSON inside production Environment packages;
    \item if using cloud storage, confirm synchronisation has completed before moving computers.
\end{enumerate}

% ============================================================
% CHAPTER 29
% ============================================================
\chapter{Security and Privacy}

\section{Passwords}

Application and Environment passwords are separate. Password verification uses salted password hashes rather than readable plaintext storage.

\section{Security questions}

Recovery answers are stored only as salted hashes. Administrators cannot view other users' recovery answers.

\section{External calendar secrets}

Secret iCal URLs use macOS Keychain integration rather than ordinary Environment data. The Secret value is machine-local; it is not transferred merely because a \fileext{.testudoenv} package is synchronized to another Mac. Since Testudo 0.1.9, the credential for an existing subscription can be restored without changing that Calendar's identity.

\section{Sensitive portable files}

A \fileext{.testudoenv} package may contain substantial professional context, including entities, relationships and history. Treat it as potentially sensitive when copying, sharing or backing it up.

% ============================================================
% CHAPTER 30
% ============================================================
\chapter{Deletion Semantics}

Deletion is permanent and differs from lifecycle transitions.

\begin{description}
    \item[Task / Note / Activity] Deleting a Work item removes it and descendant Work, together with associated relationships and history for the deleted items.
    \item[Theme] Deleting a Theme can also delete descendant Themes and Work assigned inside that Theme hierarchy according to the confirmation shown by the application.
    \item[Organization / Group] Deleting the entity removes structural memberships and Work relationships; contained entities themselves are retained where the application confirmation states so.
    \item[Person] Deletion removes the Person profile, structural memberships and Work relationships, but is blocked while the Person is linked to an Environment membership.
    \item[Calendar Event] Deletion permanently removes the event from Testudo; an externally synchronised deletion can be queued for synchronisation where supported.
\end{description}

\manualwarning{Read the confirmation text}{
Deletion scope differs by object type. Testudo's confirmation message is the authoritative description of what will be removed. For Tasks with meaningful history, prefer Completed, Discontinued and Archive instead of deletion.
}

% ============================================================
% CHAPTER 31
% ============================================================
\chapter{Recommended Working Practices}

\section{Separate plan, context and evidence}

A useful Testudo pattern is:
\begin{itemize}
    \item Task for what must be done;
    \item Note for stable context or reference material;
    \item Activity for what actually happened;
    \item Calendar Event for reserved or scheduled time.
\end{itemize}

\section{Use Discontinued instead of pretending success}

If work ended without successful completion, use Discontinued with a reason and optional outcome. This gives later readers a truthful explanation.

\section{Close only when the top-level work is no longer part of active workflow}

A Completed or Discontinued Task can remain visible in active status lists for as long as it is still operationally relevant. Close it when the whole top-level Task tree belongs in Archive.

\section{Use Themes for concepts, not as a duplicate Task tree}

Themes are semantic classification. Parent Task relationships should represent Work decomposition, while Theme hierarchy should represent conceptual grouping.

\section{Model recurring professional context explicitly}

If the same Person, lab, institution or committee repeatedly appears, create an entity and relationship rather than repeatedly typing names into free text.

\section{Record outcomes as Activities}

For long-running work, short Activities documenting important events provide a more useful record than continuously rewriting the Task description.

% ============================================================
% CHAPTER 32
% ============================================================
\chapter{Troubleshooting}

\section{The installer says the source checkout has local changes}

The automatic installer deliberately refuses to overwrite local work. Inspect the checkout:
\begin{lstlisting}
cd ~/Testudo
git status
\end{lstlisting}
Commit, stash or intentionally discard the local changes before using the automatic updater.

\section{Apple developer tools are missing}

Request Apple's Command Line Tools:
\begin{lstlisting}
xcode-select --install
\end{lstlisting}
Then rerun the compatibility check or installer.

\section{An Environment package was moved}

If the registered path or security-scoped bookmark no longer points to the package, use \button{Open Existing Environment...} and select the package at its new location.

\section{Environment login fails}

Check:
\begin{itemize}
    \item the correct Environment is active;
    \item the username is correct;
    \item the membership is Active;
    \item a password is actually configured for the membership;
    \item password recovery questions are configured if using recovery.
\end{itemize}
If recovery is unavailable, contact an Environment Administrator.

\section{Apple Calendar does not appear}

Check macOS Calendar permission, verify that the desired calendar is available in Apple Calendar, reconnect if necessary and run external-calendar synchronisation.

\section{iCal synchronisation fails}

Possible causes include an invalid or expired secret address, connectivity problems, changes at the calendar provider, or a Secret address that has not yet been stored in the current Mac's Keychain. If Testudo reports that this Mac does not have the Secret iCal address, use \button{Restore Secret Address} and enter the existing subscription URL. This restores the credential in place and should be preferred to disconnecting and recreating the Calendar when the existing Calendar identity and relationships must be preserved. If macOS asks for Keychain access after Testudo has been rebuilt, authorize the application and retry synchronization.

\section{A Task cannot be Closed}

Closing is available only for top-level Tasks whose status is Completed or Discontinued. Complete or Discontinue the Task first. Child Tasks are archived through their closed top-level ancestor.

\section{A Discontinued status change appears unavailable in a shortcut}

Discontinued requires a reason and date. Use the full Work editor so the required metadata can be entered and validated.

\section{Archived Work is missing from active lists}

This is expected. A Closed top-level Task and its subtree are removed from active Task, Today and related active-work queries. Open Archive or Reopen the root Task.

% ============================================================
% CHAPTER 33
% ============================================================
\chapter{Guided End-to-End Verification Lab}

This chapter provides a complete hands-on test that exercises the most important 0.1.7 behaviours in one disposable Environment.

\section{Step 1: create the Environment}

Create a new Environment named \emph{Manual Verification}. Create its Administrator account and enter the Environment.

Expected result: a fresh Environment opens with no Work data other than any application defaults.

\section{Step 2: create professional context}

Create:
\begin{itemize}
    \item Organization: \emph{Example Institute};
    \item Group: \emph{Field Team};
    \item Person: \emph{Alex Example};
    \item Theme: \emph{Campaign 2027}.
\end{itemize}
Affiliate the Person with the Group and/or Organization as appropriate.

Expected result: each object has a read-only detail inspector and whole-object Edit action, and navigation links work between structural objects.

\section{Step 3: create a top-level Task}

Create Task \emph{Prepare campaign} assigned to Theme \emph{Campaign 2027}. Add a relationship \term{With} Alex Example and enable inheritance for child items. Set a future Deadline.

Expected result: the Task appears in To Do and in Related Work for the selected Theme/Person where appropriate.

\section{Step 4: create children}

Inside the Task create:
\begin{itemize}
    \item sub-task \emph{Check instruments};
    \item Note \emph{Shipping instructions};
    \item Activity \emph{Kick-off discussion} with an Occurred time.
\end{itemize}

Expected result: the Task's Contains section shows all three children and inherited relationships behave according to the relationship settings.

\section{Step 5: exercise status and Today}

Change \emph{Check instruments} to In Progress. Confirm Started is populated or set manually. Create another Task due today.

Expected result: Today shows relevant active work in the appropriate Needs Attention / In Progress / Upcoming areas.

\section{Step 6: exercise Discontinued}

Create Task \emph{Use legacy sensor}. Edit it to Discontinued, reason \emph{Superseded}, outcome \emph{New sensor selected}.

Expected result: the Task appears in Discontinued, its detail page shows the metadata, and the Log contains the lifecycle event.

\section{Step 7: exercise Completed and Archive}

Complete the top-level \emph{Prepare campaign} Task, then choose Close.

Expected result:
\begin{itemize}
    \item the complete subtree moves to Archive;
    \item the root still reports Status = Completed;
    \item active Task lists no longer include the tree;
    \item Close is present in history.
\end{itemize}

\section{Step 8: navigation history}

While viewing the archived Task, open a Person, then a Theme. Use Back twice and Forward twice.

Expected result: sidebar section, middle-column selection and right-pane content travel together through the history.

\section{Step 9: Reopen}

Return to the archived root and choose Reopen.

Expected result: the Task tree leaves Archive, returns to active views with Status = Completed, and the history contains both Closed and Reopened.

\section{Step 10: calendar}

Create a local Calendar Event, link it to the Task and Theme, and set a location. Also create a Note reminder and Task deadline on the same date.

Expected result: Calendar distinguishes Event, Reminder and Deadline while still allowing navigation to the related Work.

\manualtip{Verification complete}{
If all ten steps behave as described, you have exercised the core 0.1.7 data model, Task lifecycle, Archive semantics, three-column navigation, semantic relationships, Today, Calendar and persistence-facing workflows.
}

% ============================================================
% CHAPTER 34
% ============================================================
\chapter{Technical Reference}

\section{Important file extensions}

\begin{longtable}{L{0.28\textwidth} L{0.62\textwidth}}
\toprule
\textbf{Extension} & \textbf{Purpose} \\
\midrule
\endhead
\fileext{.testudoenv} & Portable Work Environment package. \\
\fileext{.testudouser} & Portable local Testudo profile information; does not contain Environment registry or passwords. \\
\fileext{.dmg} & Locally generated macOS installer disk image. \\
\bottomrule
\end{longtable}

\section{Source tree}

\begin{lstlisting}
Sources/Testudo/
|-- Auth/
|-- Demo/
|-- Models/
|-- Resources/
|-- Store/
|-- Time/
|-- Views/
`-- TestudoApp.swift
\end{lstlisting}

\section{Build scripts}

\begin{lstlisting}
Scripts/build-app.sh
Scripts/build-release.sh
Scripts/build-local-dmg.sh
Scripts/build-manual.sh
Scripts/version.sh
\end{lstlisting}

\section{Version identity}

\filepath{Scripts/version.sh} is the build scripts' central version source. For this manual:
\begin{lstlisting}
TESTUDO_VERSION="0.1.10"
TESTUDO_BUILD_NUMBER="11"
\end{lstlisting}

\section{Task status persistence compatibility}

The current user-visible Task status set excludes Closed. A legacy persistence-only Closed value remains decodable so Environment files written by Testudo 0.1.6 and earlier can be migrated to the separate \texttt{closedAt} archive lifecycle.

% ============================================================
% CHAPTER 35
% ============================================================
\chapter{Glossary}

\begin{description}
\item[Activity] A Work item recording something that happened, optionally with an Occurred timestamp.
\item[Archive] The view containing top-level Tasks that have been Closed, together with their retained subtrees.
\item[Calendar Event] A scheduled calendar-domain object, distinct from Activity.
\item[Closed] An archive lifecycle state for terminal top-level Tasks; not a user-facing Task workflow status.
\item[Completed] A terminal Task status representing successful completion.
\item[Discontinued] A terminal Task status representing work that ended without successful completion and carries a reason/date.
\item[Entity] A Person, Group or Organization.
\item[Environment Administrator] A membership with administrative privileges inside one Work Environment.
\item[Environment membership] An identity belonging to one Work Environment, separate from the local Testudo profile.
\item[Group] A laboratory, team, committee, collaboration or similar collective entity.
\item[Local Testudo profile] The installation-level professional profile of the Testudo user.
\item[Note] An informational Work item without Task status.
\item[Organization] An institutional entity that can participate in structural hierarchy.
\item[Person] An individual represented within a Work Environment.
\item[Related Work] Work associated with a Theme or professional entity directly or through supported descendant aggregation.
\item[Reopened] A lifecycle event that returns a Closed top-level Task from Archive while preserving its real workflow status.
\item[Task] An actionable Work item with one of the user-facing statuses To Do, In Progress, Completed or Discontinued.
\item[Theme] A hierarchical semantic classification independent of Work and entity hierarchies.
\item[Work Environment] An independent portable workspace represented by a \fileext{.testudoenv} package.
\item[Work relationship] A semantic link between Work and a Person, Group or Organization.
\end{description}

% ============================================================
% APPENDICES
% ============================================================
\appendix

\chapter{Command Reference}

\section{Check compatibility}
\begin{lstlisting}
curl -fsSL \
  https://raw.githubusercontent.com/foivoskar/Testudo/main/install.sh \
  | bash -s -- --check
\end{lstlisting}

\section{Install}
\begin{lstlisting}
curl -fsSL \
  https://raw.githubusercontent.com/foivoskar/Testudo/main/install.sh \
  | bash
\end{lstlisting}

\section{Clone manually}
\begin{lstlisting}
git clone https://github.com/foivoskar/Testudo.git
cd Testudo
\end{lstlisting}

\section{Update source}
\begin{lstlisting}
git pull --ff-only
\end{lstlisting}

\section{Build local installer}
\begin{lstlisting}
./Scripts/build-local-dmg.sh
\end{lstlisting}

\section{Build application only}
\begin{lstlisting}
./Scripts/build-app.sh
\end{lstlisting}

\section{Request Apple Command Line Tools}
\begin{lstlisting}
xcode-select --install
\end{lstlisting}

\chapter{Manual Source and Compilation}

The manual source lives under:
\begin{lstlisting}
Documentation/TestudoManual/
\end{lstlisting}

The maintained source files are:
\begin{lstlisting}
Testudo_manual.tex
testudomanual.cls
\end{lstlisting}

This edition deliberately keeps \filepath{testudomanual.cls} unchanged.

The normal repository-level build command is:
\begin{lstlisting}
./Scripts/build-manual.sh
\end{lstlisting}

The script compiles the manual with XeLaTeX using the existing class file and writes the finished PDF at the repository root. A direct build remains possible:
\begin{lstlisting}
cd Documentation/TestudoManual
latexmk -xelatex Testudo_manual.tex
\end{lstlisting}

The repository-level final copy is named:
\begin{lstlisting}
Testudo_manual.pdf
\end{lstlisting}

\chapter{Quick Acceptance Checklist}

Use this checklist after installing a new Testudo build:
\begin{enumerate}
    \item Local login works and Settings are accessible.
    \item A Work Environment can be opened or created.
    \item New Entry creates Task, Note and Activity correctly.
    \item Task status filters show To Do, In Progress, Completed and Discontinued.
    \item Discontinued requires a reason and preserves outcome metadata.
    \item Only terminal top-level Tasks expose Close.
    \item Archive preserves the actual Completed/Discontinued status.
    \item Reopen returns the subtree to active views.
    \item Back/Forward restores all three columns, including Archive destinations.
    \item Today reflects current-day Work lifecycle activity.
    \item Calendar shows Events, deadlines and reminders.
    \item External calendar events remain read-only except for Testudo relationships.
    \item Person/Theme/Organization/Group Related Work navigation works.
    \item A disposable Environment can be closed and reopened without data loss.
\end{enumerate}

\chapter{Early Release Notice}

Testudo is under active development. Screens, labels and individual workflows may evolve between early versions. Before upgrading important data, keep an independent copy of critical \fileext{.testudoenv} packages.

\chapter*{License}
\addcontentsline{toc}{chapter}{License}
Testudo is distributed under the MIT License. See the repository's \filepath{LICENSE} file for the complete text.

\chapter*{Author}
\addcontentsline{toc}{chapter}{Author}
\textbf{Foivos Karakostas}

\vfill
\begin{center}
\includegraphics[width=0.09\textwidth]{\TestudoIconPath}
\vspace{4mm}

{\small Testudo User Manual\\Manual version \ManualVersion\\\ManualDate}
\end{center}

\end{document}
